\chapter{GIỚI THIỆU}

Trong bối cảnh hiện đại, dữ liệu đã trở thành một tài sản quan trọng đối với mọi tổ chức, doanh nghiệp và lĩnh vực nghiên cứu. Phân tích dữ liệu (Data Analysis) không chỉ giúp hiểu rõ hơn về tình hình hiện tại mà còn hỗ trợ ra quyết định, dự đoán xu hướng tương lai, và phát hiện các mẫu tiềm ẩn trong dữ liệu. Từ đó, phân tích dữ liệu đóng vai trò thiết yếu trong các hệ thống thông minh, mô hình dự báo, và ứng dụng học máy.

Quá trình phân tích dữ liệu thường trải qua các bước chính như sau:

\begin{enumerate}[leftmargin=1.8em,label=\textbf{(\arabic*)}]
    \item \textbf{Thu thập dữ liệu}: từ các nguồn có sẵn (file, cơ sở dữ liệu, API, cảm biến, v.v.).
    \item \textbf{Tiền xử lý dữ liệu}: làm sạch, chuyển đổi định dạng, xử lý giá trị thiếu và ngoại lệ.
    \item \textbf{Khám phá và phân tích dữ liệu}: sử dụng thống kê mô tả và trực quan hóa để hiểu cấu trúc và đặc điểm của dữ liệu.
    \item \textbf{Áp dụng mô hình}: lựa chọn và triển khai các thuật toán học máy hoặc thống kê để rút trích thông tin hoặc đưa ra dự đoán.
    \item \textbf{Đánh giá và kết luận}: kiểm tra kết quả, so sánh với các nghiên cứu khác và đưa ra nhận định.
\end{enumerate}

Dữ liệu trong thực tế tồn tại dưới nhiều dạng khác nhau. Trong phạm vi báo cáo này, chúng tôi tập trung phân tích và so sánh ba nhóm dữ liệu phổ biến:

\begin{itemize}[leftmargin=1.6em,label=\textbullet]
    \item \textbf{Nhóm A -- Dữ liệu dạng bảng (Tabular Data)}: là dạng dữ liệu có cấu trúc, thường được lưu dưới dạng bảng gồm các dòng (bản ghi) và cột (thuộc tính). Đây là loại dữ liệu thường gặp trong các hệ thống quản lý như bệnh viện, chăm sóc sức khỏe, ngân hàng. Nhóm này được đại diện bởi ba bộ dữ liệu: \textit{Heart Disease} (dự đoán bệnh tim), \textit{CDC Diabetes Health Indicators} (dấu hiệu bệnh tiểu đường) và \textit{Breast Cancer Wisconsin Diagnostic} (chẩn đoán ung thư vú).
    \item \textbf{Nhóm B -- Dữ liệu chuỗi thời gian và cảm biến (Time Series / Sensor Data)}: là dữ liệu được ghi nhận theo thời gian hoặc từ cảm biến, mỗi điểm dữ liệu tương ứng với một thời điểm cụ thể. Các ứng dụng điển hình bao gồm giám sát tiêu thụ năng lượng, nhận dạng hoạt động con người và phân tích tín hiệu sinh học. Nhóm này gồm \textit{Appliances Energy Prediction} (tiêu thụ điện), \textit{Human Activity Recognition (HAR)} (nhận dạng hoạt động) và \textit{MHEALTH} (tín hiệu sinh lý khi vận động).
    \item \textbf{Nhóm C -- Dữ liệu đa phương tiện (Multimedia Data)}: là dữ liệu phi cấu trúc như ảnh và văn bản. Phân tích nhóm này thường gắn liền với các kỹ thuật trích chọn đặc trưng (feature extraction), thị giác máy tính (Computer Vision) và xử lý ngôn ngữ tự nhiên. Nhóm này gồm \textit{HAM10000} (ảnh tổn thương da) và \textit{Reuters-21578} (văn bản phân loại tin tức).
\end{itemize}

Việc phân tích và so sánh ba nhóm dữ liệu trên sẽ giúp làm rõ điểm mạnh, điểm yếu và kỹ thuật phù hợp cho từng loại, từ đó nâng cao hiểu biết và khả năng ứng dụng vào các bài toán thực tiễn. Để chuẩn bị cho các phân tích này, Chương 2 trình bày cơ sở lý thuyết về thống kê mô tả, trực quan hóa và học máy; Chương 3 mô tả và đánh giá chất lượng của các bộ dữ liệu; Chương 4 trình bày kết quả phân tích đặc điểm dữ liệu qua thống kê và trực quan hóa; và Chương 5 ứng dụng các mô hình học máy cùng đối sánh kết quả trên cả ba nhóm dữ liệu.

